\section{Canonical one-sided histories and the connection root}
\label{sec:connection}

The reduced RFDE dynamics on the history graph is a complete planar flow,
but the singular canard still has two linear obstructions: translation
along the orbit and a normally growing Gaussian mode.  We remove them by
one common phase condition and opposite one-sided normal conditions.  This
constructs the preparation-indexed histories used in
\cref{thm:main} and turns the return coefficient
\cref{eq:transverse-return-coefficient} into a root displacement.

\subsection{An exact tangent--normal Green operator}

Linearization of \(q_0\) along \(\gamma_0\) gives
\begin{equation}
\label{eq:L0}
 L_0\xi=f,\qquad
 L_0\vect{U\\V}=\vect{U'-sU-V\\V'+U}.
\end{equation}
\begin{samepage}
Put
\begin{equation}
\label{eq:fundamental-modes}
 E(s)=\e^{s^2/2},\qquad I(s)=\int_0^s\e^{t^2/2}\,\dd t,
\end{equation}
and
\begin{equation}
\label{eq:tau-n}
 \tau(s)=\vect{-1\\s},
 \qquad n(s)=\vect{I(s)\\E(s)-sI(s)}.
\end{equation}
\end{samepage}
Both are homogeneous solutions of \cref{eq:L0}; \(\tau=2\alpha\gamma_0'\)
is tangent and \(n\) is normal.  Since
\(\det[\tau,n]=-E\), every solution can be written uniquely as
\(\xi=a\tau+bn\), with
\begin{equation}
\label{eq:ab-coordinates}
\begin{aligned}
 a&=E^{-1}\bigl[-(E-sI)U+IV\bigr],\\
 b&=E^{-1}(sU+V).
\end{aligned}
\end{equation}
Variation of constants becomes
\begin{equation}
\label{eq:ab-variation}
 a'=E^{-1}\bigl[-(E-sI)f_1+If_2\bigr]=:A_f,
 \qquad
 b'=E^{-1}(sf_1+f_2)=:B_f.
\end{equation}

On \(J_a=[-R,0]\) and \(J_r=[0,R]\), impose the common phase
\(X(0)=0\).  At the linear level this is exactly \(a(0)=0\).  Impose
separately \(b(-R)=\beta_a\) on the left and \(b(R)=\beta_r\) on the right.
The corresponding Green operators are
\begin{equation}
\label{eq:one-sided-green}
\begin{aligned}
 a_a(s)&=-\int_s^0A_f(t)\,\dd t,&
 b_a(s)&=\beta_a+\int_{-R}^sB_f(t)\,\dd t,\\
 a_r(s)&= \int_0^sA_f(t)\,\dd t,&
 b_r(s)&=\beta_r-\int_s^RB_f(t)\,\dd t.
\end{aligned}
\end{equation}

For \(c\ge0\) and \(m\in\mathbb N_0\), let
\begin{equation}
\label{eq:tame-norms}
 \norm{g}_{c,m,J}
 =\sup_{s\in J}\e^{-c|s|}\langle s\rangle^{-m}|g(s)|,
 \qquad
 |\beta|_{c,m,R}
 =\e^{R^2/2-cR}\langle R\rangle^{-m}|\beta|.
\end{equation}

\begin{lemma}[Uniform phase--normal Green estimate]
\label{lem:green}
For fixed \(c,m\), there are \(C>0\) and \(d_0\le4\), independent of
\(R\ge1\), such that
\begin{equation}
\label{eq:green-bound}
 \norm{\cG_R^\sigma(f,\beta_\sigma)}_{c,m+d_0,J_\sigma}
 \le C\bigl(\norm{f}_{c,m,J_\sigma}+|\beta_\sigma|_{c,m,R}\bigr),
 \qquad \sigma=a,r.
\end{equation}
The estimate holds componentwise for the finite rectangular family of
\((\nu,\eta)\)-derivatives used below.
\end{lemma}

\begin{proof}
Integration by parts gives
\begin{equation}
\label{eq:I-bounds}
 |I(s)|\le CE(s)\langle s\rangle^{-1},
 \qquad |E(s)-sI(s)|\le CE(s).
\end{equation}
Thus \(|A_f|\le C|f|\) and
\(|B_f|\le CE^{-1}\langle s\rangle|f|\).  The tangent integrals in
\cref{eq:one-sided-green} cost only polynomial weights.  For the normal
integrals, splitting at \(|s|=1\) and integrating once more by parts gives
\begin{align}
 E(s)\int_{-R}^sE(t)^{-1}\e^{c|t|}\langle t\rangle^{m+1}\,\dd t
 &\le C\e^{c|s|}\langle s\rangle^{m+2}, &&-R\le s\le0,
 \label{eq:left-gaussian-tail}\\
 E(s)\int_s^RE(t)^{-1}\e^{c|t|}\langle t\rangle^{m+1}\,\dd t
 &\le C\e^{c|s|}\langle s\rangle^{m+2}, &&0\le s\le R.
 \label{eq:right-gaussian-tail}
\end{align}
The normalized boundary factor in \cref{eq:tame-norms} cancels the largest
value of \(E\).  Taking \(d_0=4\) also covers the compact middle interval.
The kernels are parameter independent, so the same argument applies after
the declared differentiations.
\end{proof}

The one-sided placement in \cref{eq:one-sided-green} is essential.  A
two-sided unweighted fundamental-matrix estimate would introduce
\(\e^{R^2/2}\), precisely the loss that would destroy the logarithmic
matching argument.

\subsection{Prepared traces and their exact RFDE lift}

Define the first integral of the singular planar problem by
\begin{equation}
\label{eq:first-integral}
 \mathscr H(X,Y)=\frac12\e^{-2\alpha Y}
 \left(\alpha Y-\alpha^2X^2+\frac12\right).
\end{equation}
It is constant along \(q_0\), and
\begin{equation}
\label{eq:first-integral-gradient}
 \nabla\mathscr H(\gamma_0(s))
 =\frac{\alpha\e}{2}\e^{-s^2/2}(s,1)^T.
\end{equation}
Write \(S=S_\del\) and \(R=R_\del=S+B\).  The prepared reduced field equals the actual graph
field on the shrinking retained tube and equals \(q_0\) near
\(\gamma_0(\pm R)\).  It is the unique field defined by
\cref{eq:prepared-planar-field}.  The one-sided problems are
\begin{equation}
\label{eq:trace-bvp}
\begin{aligned}
 z'&=Q^{\rm pr}_{\del,\nu,\eta}(z),\\
 X^a(0)&=X^r(0)=0,\\
 \mathscr H(z^a(-R))&=0,
 &\mathscr H(z^r(R))&=0.
\end{aligned}
\end{equation}

For the exact retained pieces, use the transition sections
\begin{equation}
\label{eq:transition-sections}
 \Sigma_S^a:\quad
 Y=\frac{S^2-2}{4\alpha},\ X>0,
 \qquad
 \Sigma_S^r:\quad
 Y=\frac{S^2-2}{4\alpha},\ X<0.
\end{equation}
Let \(t_a<0<t_r\) be their unique hits by \(z^a\) and \(z^r\),
respectively, and define
\begin{equation}
\label{eq:retained-intervals}
 I_a=[t_a,0],\qquad I_r=[0,t_r].
\end{equation}
Existence, uniqueness, and the parameter estimates for these moving hits
are proved in \cref{app:moving-hits}.
The endpoint condition is exactly the nonlinear continuation of the
no-incoming-normal condition.  Indeed, if
\(z=\gamma_0+(U,V)=\gamma_0+a\tau+bn\) at an endpoint \(s\), then
\begin{equation}
\label{eq:nonlinear-boundary}
 E(s)b=\alpha[-a+I(s)b]^2.
\end{equation}

\begin{theorem}[Canonical one-sided traces]
\label{thm:canonical-traces}
For the fixed admissible preparation in \cref{def:admissible-preparation},
and after decreasing \(\del_0\), the problems
\cref{eq:trace-bvp} have unique solutions in a tame neighborhood of
\(\gamma_0\).  They possess the rectangular derivatives
\(C_\nu^1C_\eta^2\), including the mixed derivatives, and for fixed
\(c,m,C\), independent of \(\del\),
\begin{equation}
\label{eq:trace-estimates}
\begin{aligned}
 \norm{z^\sigma-\gamma_0}_{c,m,J_\sigma}
 +\norm{\partial_\nu z^\sigma}_{c,m,J_\sigma}&\le C\del,\\
 \norm{\partial_\eta z^\sigma}_{c,m,J_\sigma}
 +\norm{\partial_{\eta\eta}z^\sigma}_{c,m,J_\sigma}&\le C\del^2,\\
 \norm{\partial_\nu\partial_\eta z^\sigma}_{c,m,J_\sigma}
 +\norm{\partial_\nu\partial_{\eta\eta}z^\sigma}_{c,m,J_\sigma}
 &\le C\del^2,\qquad \sigma=a,r.
\end{aligned}
\end{equation}
By the fixed enlarged target in \cref{def:admissible-preparation}, the
current points on \(I_a\cup I_r\) lie in
\(\abs{\chi}\le\widehat S_\del\) and in the shrinking
normal tube.  All continuous delay backtracks used by
\cref{eq:history-embedding} lie in the still larger nested uncut flow hull
from \cref{def:admissible-preparation}.  Hence
\begin{equation}
\label{eq:lifted-traces}
 \varphi^\sigma(s)=\iota_{\del,\nu,\eta}(z^\sigma(s))
\end{equation}
are exact complete-history solutions of \cref{eq:physical-model} on
\(I_a\) and \(I_r\), respectively.  The right-hand trace is backward
extendible there
inside the finite-dimensional history graph.
\end{theorem}

\begin{proof}
Write \(z=\gamma_0+\xi\).  The equation is
\(L_0\xi=N_{\del,\nu,\eta}(s,\xi)\), where
\begin{equation}
\label{eq:N-bounds}
 |N(s,0)|\le C\del\e^{c_0|s|}\langle s\rangle^{m_0},
 \qquad
 |D_\xi N(s,\xi)|\le
 C\bigl(\del\e^{c_0|s|}\langle s\rangle^{m_0}+|\xi|\bigr).
\end{equation}
Equation \cref{eq:nonlinear-boundary} defines
\(b=\mathcal B(a)\), with
\(\mathcal B(0)=D\mathcal B(0)=0\), in the normalized boundary norm.
Choose weights \(c>c_0\) and \(M>m_0+d_0+2\).  Combining
\cref{lem:green,eq:N-bounds} gives a self-map of the ball
\(\norm{\xi}_{c,M}\le C_0\del\), whose Lipschitz factor is
\begin{equation}
\label{eq:trace-contraction-factor}
 \rho_{\del,S}=C\del\e^{c(S+B)}\langle S+B\rangle^M=o(1).
\end{equation}
The boundary graph has the same factor because its quadratic term contains
\(E(R)^{-1}\).  Banach's theorem proves existence and uniqueness.

Differentiating this fixed-point equation is legitimate for the finite
parameter family furnished by \cref{thm:history-graph}.  The \(\nu\)-source
is \(O(\del)\), the \(\eta\)-source is \(O(\del^2)\), and a second
\(\eta\)-derivative contains only an \(O(\del^2)\) direct source plus
products of already bounded lower derivatives.  The inverse of the
linearized contraction has a uniform bound, proving
\cref{eq:trace-estimates}; the detailed estimates are proved in
\cref{app:green-gap}.

In the canard coordinate \(d\), the exact identity
\begin{equation}
\label{eq:normal-bootstrap}
 d(\gamma_0(s)+(U,V))=V+sU-\alpha U^2
\end{equation}
and \cref{eq:trace-estimates} give
\(|d|\le C\del\e^{cS}\langle S\rangle^{m+1}
=o(\del^{\kappa_0})\).  The same nested estimate holds for the depth-two
backtracks.  Thus the retained trace never uses the artificial normal
extension.  The moving-hit estimate in \cref{eq:hitting-time-bounds} and
the construction at target \(\widehat S_\del\) place its current points in the
enlarged uncut region; the nested construction places their delay
backtracks in its outer flow hull.  Equations
\cref{eq:history-embedding,eq:exact-semiconjugacy} then give the exact
history lift \cref{eq:lifted-traces}; no backward RFDE initial-value problem
is solved.
\end{proof}

\subsection{The normalized history gap}

For the fixed preparation \(\cP\), define
\begin{equation}
\label{eq:canonical-gap}
 \mathfrak d_\cP(\del,\nu,\eta)
 =\frac{2}{\alpha\e}
 \left[\mathscr H(z^a_\cP(0))-\mathscr H(z^r_\cP(0))\right].
\end{equation}

\begin{lemma}[Gap zero is complete-history equality]
\label{lem:gap-history}
For sufficiently small \(\del\),
\begin{equation}
\label{eq:gap-history-equivalence}
 \mathfrak d_\cP=0
 \quad\Longleftrightarrow\quad
 z^a_\cP(0)=z^r_\cP(0)
 \quad\Longleftrightarrow\quad
 \iota(z^a_\cP(0))=\iota(z^r_\cP(0)).
\end{equation}
\end{lemma}

\begin{proof}
Both central points satisfy \(X=0\).  If
\(h(Y)=\mathscr H(0,Y)\) and \(Y_0=-1/(2\alpha)\), then
\(h'(Y_0)=\alpha\e/2\ne0\).  The trace theorem places the two values in a
common interval on which \(h\) is injective.  This proves equality of the
current reduced states.  Present-time evaluation followed by projection to
\(u\) is a left inverse of \(\iota\), proving the history equivalence.
\end{proof}

\begin{proposition}[Uniform gap expansion]
\label{prop:gap-expansion}
Let
\begin{equation}
\label{eq:nu0}
 \nu_0=-\frac{11}{24\alpha},
 \qquad
 c_\perp=\frac{K(\theta_0-\theta_1)}{4\alpha}.
\end{equation}
Uniformly for \(\nu\) in a fixed neighborhood of \(\nu_0\) and
\(|\eta|\le\eta_0\),
\begin{equation}
\label{eq:full-gap-expansion}
\begin{aligned}
 \frac{\mathfrak d_\cP}{\del}
 &=\sqrt{2\pi}(\nu-\nu_0)+O(\del),\\
 \partial_\nu\mathfrak d_\cP
 &=\sqrt{2\pi}\,\del+O(\del^2),\\
 \partial_\eta\mathfrak d_\cP
 &=-c_\perp\sqrt{2\pi}\,\del^2
   +O(\del^3+\del^2|\eta|),\\
 \partial_{\eta\eta}\mathfrak d_\cP&=O(\del^2).
\end{aligned}
\end{equation}
\end{proposition}

\begin{proof}
Because both prepared endpoints lie on \(\mathscr H=0\), the fundamental
theorem of calculus gives the exact two-piece identity
\begin{equation}
\label{eq:two-piece-gap}
 \frac{\alpha\e}{2}\mathfrak d_\cP
 =\int_{-R}^{0}\nabla\mathscr H(z^a)^T
       (Q^{\rm pr}-q_0)(z^a)\,\dd s
  +\int_{0}^{R}\nabla\mathscr H(z^r)^T
       (Q^{\rm pr}-q_0)(z^r)\,\dd s.
\end{equation}
The prepared tails are supported outside the retained core; the Green
formula weights their central contribution by
\(O(\e^{-S^2/2+cS}\langle S\rangle^m)\).  After choosing \(p>p_*\), this
and all differentiated endpoint terms are smaller than the remainders in
\cref{eq:full-gap-expansion}.  On the retained core,
\cref{eq:growing-graph-expansion,eq:trace-estimates} and the Gaussian
gradient \cref{eq:first-integral-gradient} provide a common integrable
majorant.  Hence the traces may be replaced by \(\gamma_0\) at the stated
cost, and the finite integrals may be extended to \(\R\).

The first graph coefficient on \(\gamma_0\) is
\begin{equation}
\label{eq:Q1-on-canard}
 Q_{1,X}(\gamma_0(s))
 =\frac{11s^3-12K(\theta_0+2\theta_1)}{72\alpha},
 \qquad Q_{1,Y}(\gamma_0(s))=\nu.
\end{equation}
Pairing with \(\e^{-s^2/2}(s,1)^T\), the constant delay term is odd,
whereas
\[
 \int_\R\e^{-s^2/2}\,\dd s=\sqrt{2\pi},
 \qquad
 \int_\R s^4\e^{-s^2/2}\,\dd s=3\sqrt{2\pi}.
\]
This gives the first line of \cref{eq:full-gap-expansion}.  Differentiating
\cref{eq:two-piece-gap} with respect to \(\nu\) leaves at order \(\del\)
only
\(\del\nabla\mathscr H(\gamma_0)^T(0,1)\), proving the second line.

Since \(Q_1\) is independent of \(\eta\), its differentiated trace term is
\(O(\del^3)\).  The unique order-\(\del^2\) source is
\(
 \del^2\nabla\mathscr H(\gamma_0)^T\partial_\eta Q_2(\gamma_0)
\).
Using \cref{eq:transverse-return-coefficient}, its normalized pairing is
\begin{equation}
 -c_\perp\del^2\int_\R s^2\e^{-s^2/2}\,\dd s
 =-c_\perp\sqrt{2\pi}\,\del^2.
\end{equation}
The trace replacement, \(R_Q\), and variation away from \(\eta=0\) give
the three error terms displayed in the third line.  A second
\(\eta\)-derivative contains an \(O(\del^2)\) direct source; all terms
involving \(\partial_\eta z\) or
\(\partial_{\eta\eta}z\) are of that order or smaller by
\cref{eq:trace-estimates}.  This proves the final line.  Detailed endpoint
and moving-hit estimates are in \cref{app:green-gap}.
\end{proof}

\subsection{Proof of the main theorem}

\begin{proof}[Proof of \cref{thm:main}]
Choose a compact \(\nu\)-interval centered at \(\nu_0\).  By
\cref{prop:gap-expansion}, after reducing \(\del_0\),
\begin{equation}
\label{eq:gap-monotonicity}
 \partial_\nu\mathfrak d_\cP
 \ge\frac{\sqrt{2\pi}}2\del
\end{equation}
throughout a fixed neighborhood
\(|\nu-\nu_0|\le c_0\).  The first line of
\cref{eq:full-gap-expansion} gives opposite signs at its endpoints for
small \(\del\).  Hence there is exactly one root
\(\nu_{c,\cP}(\del,\eta)\) in that interval.  For every fixed
\(0<\del\le\del_0\), the ordinary implicit-function theorem and
\cref{eq:gap-monotonicity} make the root \(C^1\) in \(\eta\); no joint
smoothness in \((\del,\eta)\) is asserted because the target-dependent
cutoff was frozen before expanding in the dummy amplitude.

The first line also gives \(\nu_{c,\cP}=\nu_0+O(\del)\), uniformly in
\(\eta\).  Implicit differentiation and the remaining lines give
\begin{equation}
 \partial_\eta\nu_{c,\cP}
 =-\frac{\partial_\eta\mathfrak d_\cP}
         {\partial_\nu\mathfrak d_\cP}
 =c_\perp\del+O(\del^2+\del|\eta|).
\end{equation}
Integrating from \(0\) to \(\eta\) and using
\(\mu_{c,\cP}=\del^2\nu_{c,\cP}\) proves
\cref{eq:main-law} with the claimed uniform remainder.

At the root, \cref{lem:gap-history} identifies the two central complete
histories.  Uniqueness for the autonomous planar \(Q\)-flow concatenates
the two retained pieces into one local reduced orbit, and the exact
semiconjugacy \cref{eq:exact-semiconjugacy} maps it to one local RFDE orbit.
\end{proof}
